\section{Cleaning arbitrary incidence systems}
\label{sec:cleaning}
The next construction is \cite[Theorem 4.2]{DMV}. We include the rigidity of
its tail and the isolated-vertex case, since those are the local obligations
that make its symmetry assertion valid.

\begin{theorem}[Cleaning]\label{UR.Cleaning}
Every nonempty finite properly colored graph $X$ has a finite properly colored,
triangle-free graph $L(X)$ of minimum degree at least two satisfying
\[
 (C(L(X)),K(L(X)))\cong(C(X),K(X)).
\]
It preserves the connected-component decomposition and adds exactly
$q(2M+5)$ vertices, where
\[
 M=\max_{i\in I}|t^{-1}(i)|,
 \qquad W=\{v:\deg_X(v)\leq1\},\qquad q=|E(X)|+|W|.
\]
Its palette consists of the old colors and two new colors.
\end{theorem}
\begin{proof}
Since $X$ is nonempty, $M\geq1$; also $q\geq1$, since otherwise a nonempty
edgeless graph would have no vertices of degree at most one.
For each edge $e$ and each $w\in W$ take a separate gadget with vertices
$u_0,\ldots,u_{2M+4}$, consecutive edges $u_ju_{j+1}$, and the two chords
\begin{equation}\label{eq:tail}
 u_{2M+4}u_{2M-1},\qquad u_{2M+3}u_{2M}.
\end{equation}
Delete all old edges. For $e=\{v,w\}$ join both endpoints to the root
$u_{e,0}$ of its edge gadget. Join each $w\in W$ to $u_{w,0}$; if $w$ was
isolated, also join it to $u_{w,2}$. Retain old colors, color even gadget
indices by $\alpha$ and odd indices by $\beta$, where $\alpha,\beta$ are
distinct new colors.

Every new edge joins different colors. An old vertex of degree at least two
retains that degree; a former leaf acquires a second neighbor; a former
isolated vertex acquires its two displayed neighbors. Every gadget vertex has
degree at least two. There are no old--old edges. The new induced graph is
bipartite. A triangle containing an old vertex would require two adjacent new
neighbors in one gadget, but the only old vertex with two neighbors in one
gadget was isolated, and its neighbors $u_0,u_2$ are not adjacent. Thus $L(X)$
is triangle-free and has the required minimum degree.

The two new color classes have sizes $q(M+3)$ and $q(M+2)$. These sizes are
distinct and both exceed $M$, the largest old class size. Every correlation
therefore fixes each of these classes and preserves the old vertex set. Two
distinct old vertices were adjacent in $X$ exactly when they have a common
new neighbor in $L(X)$: their edge gadget supplies its root, and every other
gadget has only one old neighbor. This reconstructs the old graph as well as
its color partition, so restriction belongs to $C(X)$.

To establish rigidity, delete all old vertices. The components are precisely
the separate gadgets. Each edge gadget has its old endpoint pair as boundary;
each additional gadget has its one old vertex as boundary. The graph is simple
and at most one additional gadget is attached to a given old vertex, so all
these boundaries are distinct.

Consider a correlation of a gadget together with its boundary that fixes the
boundary. For an edge gadget or a former leaf, the root is the unique gadget
vertex adjacent to the boundary. For an isolated vertex the two attached
vertices $u_0,u_2$ have different degrees in this induced graph: the first has
degree two and the second at least three. Hence the root is fixed in all cases.
Starting at the root forces successive vertices until $u_{2M-1}$. At that
vertex the forward neighbors $u_{2M}$ and $u_{2M+4}$ have induced-gadget
degrees three and two respectively, so both are fixed. At $u_{2M}$ the
remaining neighbors $u_{2M+1}$ and $u_{2M+3}$ have induced-gadget degrees two
and three respectively. Finally $u_{2M+2}$ is forced. These branch degrees
are taken within the gadget, after its component has been recognized; an
isolated attachment therefore causes no exception when $M=1$.
Lemma~\ref{UR.Components} now gives pointwise rigidity over all old vertices.

A correlation $f$ of $X$ preserves degrees. Extend it by
\[
 u_{e,j}\longmapsto u_{f(e),j},\qquad
 u_{w,j}\longmapsto u_{f(w),j}.
\]
The complete edge list, including the isolated attachment, is preserved, and
the inverse is the same construction from $f^{-1}$. Both new colors are fixed;
old colors are fixed exactly when $f$ fixes them. Thus
Theorem~\ref{UR.Restriction} gives the asserted pair isomorphism.
There are $q$ gadgets of $2M+5$ vertices. Replacing an old edge preserves
connectivity between its endpoints and does not join different old components;
all other gadgets attach to one old vertex. This proves the component claim.
\end{proof}

\begin{corollary}\label{UR.Conversion}
Every finite incidence system has a finite incidence geometry with the same
correlation group and the same distinguished type-preserving subgroup.
Connected inputs give connected outputs.
\end{corollary}
\begin{proof}
Apply Theorem~\ref{UR.Chambers} to $L(X)$ from Theorem~\ref{UR.Cleaning}.
Both restriction maps identify the distinguished kernels. Each construction
preserves the component decomposition.
\end{proof}
